\paragraph{Lee--Yang 端点坐标与共动偏移的平方根统一化}
\begin{theorem}[共动 Cayley 实像的 Lee--Yang 端点定标]\label{thm:xi-comoving-leyang-endpoint-square-root}
\leanverified{paper\_xi\_comoving\_leyang\_endpoint\_square\_root}
沿用 \eqref{eq:xi-comoving-pure-delta} 的记号，并令 \(J\) 为附录 \eqref{eq:lee-yang-J-def} 的坐标门。
则对任意离临界线零点 \(\rho=\beta+\mathrm{i}\gamma\)（\(\beta<\tfrac12\)）及其偏移 \(\delta=\tfrac12-\beta\) 有严格恒等式
\[
J(w_{\rho,\gamma})=-\frac{1-\delta^2}{4},
\qquad
\delta=\sqrt{\,1+4J(w_{\rho,\gamma})\,},
\qquad
J(w_{\rho,\gamma})+\tfrac14=\frac{\delta^2}{4}.
\]
因此 \(\delta\) 在 \(J\)-坐标中对应于端点 \(-\tfrac14\) 的平方根统一化参数。
\end{theorem}
\begin{proof}
由 \eqref{eq:xi-comoving-pure-delta} 有 \(w_{\rho,\gamma}=\frac{1-\delta}{1+\delta}\)。
计算
\[
1+w_{\rho,\gamma}=\frac{(1+\delta)+(1-\delta)}{1+\delta}=\frac{2}{1+\delta}.
\]
代入 \(J(u)=-u/(1+u)^2\) 得
\[
J(w_{\rho,\gamma})
=-\frac{\frac{1-\delta}{1+\delta}}{\left(\frac{2}{1+\delta}\right)^2}
=-\frac{(1-\delta)(1+\delta)}{4}
=-\frac{1-\delta^2}{4}.
\]
移项即得 \(1+4J(w_{\rho,\gamma})=\delta^2\)，从而 \(\delta=\sqrt{1+4J(w_{\rho,\gamma})}\)；
再作一次移项得到 \(J(w_{\rho,\gamma})+\tfrac14=\delta^2/4\)。证毕。
\end{proof}

\begin{conjecture}[Lee--Yang 端点能量谱的平方根刚性]\label{conj:xi-leyang-endpoint-square-root-rigidity}
在有限缺陷口径下，每个离临界线零点 \(\rho=\beta+\mathrm{i}\gamma\)（按重数计）在共动图表 \(x_0=\gamma\) 上诱导一个端点能量
\[
\varepsilon(\rho):=J(w_{\rho,\gamma})+\tfrac14=\frac{\delta(\rho)^2}{4}>0,
\qquad \delta(\rho):=\tfrac12-\Re(\rho).
\]
令 \(\mathscr E\) 为所有此类 \(\varepsilon(\rho)\) 的多重集。
则存在一类仅依赖边界谱数据（而不显式访问盘内零点位置）的可审计泛函 \(\mathfrak F\)，满足
\[
\mathfrak F=0 \quad\Longleftrightarrow\quad \sum_{\varepsilon\in\mathscr E}\varepsilon=0.
\]
由于 \(\varepsilon>0\)，右端等价于 \(\mathscr E=\varnothing\)，从而把 RH 压缩为单端点能量为零的谱条件。
\end{conjecture}

\begin{remark}\label{rem:xi-leyang-endpoint-square-root-rigidity-route}
该猜想要求把“离线偏移”由 \(\delta\) 重参数化为端点能量 \(\varepsilon=\delta^2/4\)，并在边界证书层面实现对 \(\varepsilon\)-谱的可逆读出。
若某个已构造的共动边界泛函能够被显式改写为关于 \(\{\varepsilon(\rho)\}\) 的可逆变换（例如在有限缺陷口径下给出 \(\sum_{\rho}m(\rho)\Phi(\varepsilon(\rho))\) 型表达且 \(\Phi\) 在相关范围上可逆），则端点能量零化与离线零点不存在性可在同一审计框架内闭合。
\end{remark}

\endinput
